% Chapter 4: From Brane Tension to Pinning Constant
% Source: M6_K_RIGOROUS_DERIVATION.md, M6_PINNING_CONSTANT_DERIVATION.md
% Status: FILLED from SoT
% Contamination check: PASSED (no SM terms)

\chapter{From Brane Tension to Pinning Constant}
\label{ch:sigma_to_K}

\source{M6\_K\_RIGOROUS\_DERIVATION.md}

\begin{abstract}
The brane tension $\bsigma$ is the fundamental energy scale of EDC. This chapter
derives the pinning constant $\Kpin$ from $\bsigma$ and contact geometry, establishing
the quantitative link between 5D brane physics and cluster binding energies.
\end{abstract}

% ----------------------------------------------------------------------------
% CHAPTER SPINE
% ----------------------------------------------------------------------------
\paragraph{Chapter Spine.}
\textbf{Purpose:} Derive the pinning constant $\Kpin$ from brane tension $\bsigma$
via contact geometry, establishing the energy scale for cluster binding.
\textbf{Inputs:} (i)~Brane tension $\bsigma = 8.82\;\MeV/\text{fm}^2$~[Dc],
(ii)~junction scale $L_0 = 1\;\text{fm}$~[P], (iii)~brane thickness $\delta = 0.105\;\text{fm}$~[P].
\textbf{Outputs:} (i)~Contact area $A_{\text{contact}} = \pi\delta L_0$~[Dc],
(ii)~geometric factor $f = \sqrt{\delta/L_0}$~[I], (iii)~$\Kpin = f \times \bsigma \times A_{\text{contact}} \approx 0.93\;\MeV$~[Dc/I].
\textbf{Dependencies:} Ch.~\ref{ch:junction} ($\Msix$ lattice structure).
\textbf{Forward links:} $\Kpin$ to Ch.~\ref{ch:deuterium} (deuterium binding),
Ch.~\ref{ch:helium4} (closed-4 energy budget), and Ch.~\ref{ch:frustrated} (frustration correction).

% ============================================================================
\section{The Central Connection}
\label{sec:sigmaK:central}
% ============================================================================

The key insight of the topological pinning model: brane tension $\bsigma$
(established in Book~I as the fundamental EDC energy scale) controls
cluster binding through the pinning constant $\Kpin$.

\begin{derivationbox}[title=Central Claim]
    All binding energies in the cluster sector trace to a single parameter:
    \begin{equation}
        \bsigma = 8.82\;\MeV/\text{fm}^2 \tagDc
    \end{equation}
    The pinning constant $\Kpin$ emerges from the contact geometry between
    neighboring junctions in the M$_6$ lattice.
\end{derivationbox}

The derivation chain:
\begin{equation}
    \bsigma \;\longrightarrow\; A_{\text{contact}} \;\longrightarrow\; f \;\longrightarrow\; \Kpin
\end{equation}

% ============================================================================
\section{Contact Geometry}
\label{sec:sigmaK:contact}
% ============================================================================

\begin{tcolorbox}[colback=gray!5,colframe=gray!50!black,title=Notation: Length Scales $\delta$ and $L_0$]
This chapter uses two geometric parameters:
\begin{itemize}
    \item $\delta \approx 0.105\;\text{fm}$: brane thickness / flux tube radius.
          \textbf{Source:} Postulated as $\delta = \hbar/(2 m_p c)$ (Compton regularization);
          analyzed in Ch.~\ref{ch:L0delta}. \tagP
    \item $L_0 \approx 1.0\;\text{fm}$: characteristic junction scale.
          \textbf{Source:} Postulated via the ratio hypothesis $L_0/\delta = \pi^2$
          in Ch.~\ref{ch:L0delta}. \tagP
\end{itemize}
Both are \emph{imported} into this chapter; their derivation status is [P] (postulated).
The ratio $L_0/\delta \approx 9.87$ drives the instanton exponent in Ch.~\ref{ch:tau_n}.
\end{tcolorbox}

\subsection{The Physical Setup}

From the M$_6$ geometry derivation (Ch.~\ref{ch:junction}):
\begin{itemize}
    \item Each Y-junction has 6 neighbors in the dual lattice
    \item Neighbors share a \textbf{contact surface} where flux tubes meet
    \item The contact surface has tension $\bsigma$
\end{itemize}

\subsection{Primal Bond Structure}

In the primal Steiner graph, each Y-junction has 3 legs:
\begin{itemize}
    \item \textbf{Radius:} $r_{\text{tube}} \approx \delta = 0.105\;\text{fm}$
    \item \textbf{Length:} $\ell_{\text{tube}} \approx L_0/3 \approx 0.33\;\text{fm}$
\end{itemize}

\subsection{Contact Surface Models}

When two junctions connect (leg-to-hub contact), three models apply:

\begin{description}
    \item[Model A: Circular disk] Area $A = \pi\delta^2 \approx 0.035\;\text{fm}^2$
    \item[Model B: Cylindrical wrap] Area $A = 2\pi\delta \ell_{\text{contact}}$
    \item[Model C: Saddle surface] Geometric mean of the above
\end{description}

\subsection{The Geometric Mean Principle}

\begin{derivationbox}[title=Contact Area Derivation]
    From energy minimization, the contact surface is saddle-shaped. The
    characteristic scale is the geometric mean:
    \begin{equation}
        r_{\text{contact}} = \sqrt{\delta \cdot L_0} = \sqrt{0.105 \times 1.0} \approx 0.32\;\text{fm} \tagDc
    \end{equation}
    The contact area:
    \begin{equation}
        A_{\text{contact}} = \pi \cdot r_{\text{contact}}^2 = \pi \cdot \delta \cdot L_0 \approx 0.33\;\text{fm}^2 \tagDc
        \label{eq:contact_area}
    \end{equation}
\end{derivationbox}

% ============================================================================
\section{Energy of Contact}
\label{sec:sigmaK:energy}
% ============================================================================

\subsection{Matched vs.\ Mismatched States}

Consider two neighboring junctions with reaction coordinates $q_i$ and $q_j$.
Define the mismatch parameter $\Delta q := q_i - q_j$.

\begin{description}
    \item[Matched states] ($\Delta q = 0$): Contact surface is flat
        \begin{equation}
            E_{\text{matched}} = \bsigma \times A_{\text{contact}}
            \label{eq:E_matched}
        \end{equation}
    \item[Mismatched states] ($\Delta q \neq 0$): Contact surface must curve
        to accommodate the angular difference between junctions.
\end{description}

\subsection{Scaling Bridge: Why Quadratic in $\Delta q$?}
\label{sec:sigmaK:scaling_bridge}

The mismatch energy depends on $(\Delta q)^2$ for small $\Delta q$. Here is the
scaling argument:

\begin{enumerate}
    \item \textbf{Symmetry constraint:} The energy must be symmetric under
          $\Delta q \to -\Delta q$ (exchanging junction labels). Therefore,
          the leading-order correction is \emph{even} in $\Delta q$.

    \item \textbf{Curvature estimate:} When $\Delta q \neq 0$, junction angles
          deviate from the equilibrium 120°. The contact surface curves with
          characteristic radius:
          \begin{equation}
              R_c \sim \frac{L_0}{|\Delta q|}
              \label{eq:R_c}
          \end{equation}
          This is a dimensional estimate \tagP: the only length scale available
          for the curvature is $L_0$, and it must diverge as $\Delta q \to 0$.

    \item \textbf{Curvature energy:} Bending a surface of area $A$ with tension
          $\sigma$ and curvature $1/R$ costs energy $\sim \sigma A / R^2$.
          Substituting $R_c \sim L_0/|\Delta q|$:
          \begin{equation}
              E_{\text{curv}} \sim \bsigma \cdot A_{\text{contact}} \cdot
              \frac{(\Delta q)^2}{(L_0/L_0)^2}
              = \bsigma \cdot A_{\text{contact}} \cdot (\Delta q)^2
          \end{equation}
          up to $O(1)$ geometric factors. \tagP
\end{enumerate}

\textbf{Summary:} The quadratic dependence on $\Delta q$ follows from symmetry
and dimensional analysis. Unknown $O(1)$ prefactors are absorbed into the
definition of $\Kpin$ below.

\subsection{Total Mismatch Energy}

\begin{derivationbox}[title=Mismatch Energy]
    Combining the matched energy \eqref{eq:E_matched} with the curvature
    correction, the mismatch energy cost is:
    \begin{equation}
        \Delta E := E_{\text{mismatched}} - E_{\text{matched}}
                  = \bsigma \times A_{\text{contact}} \times (\Delta q)^2
        \label{eq:mismatch_energy}
    \end{equation}
    up to $O(1)$ geometric factors absorbed into later definitions. \tagP

    This is the \textbf{pinning energy per bond}.
\end{derivationbox}

% ============================================================================
\section{The Pinning Formula}
\label{sec:sigmaK:formula}
% ============================================================================

\subsection{Definition of $\Kpin$}

The pinning Hamiltonian couples neighboring junctions:
\begin{equation}
    H_{\text{pin}} = \Kpin \sum_{\langle i,j \rangle} (q_i - q_j)^2
    \label{eq:H_pin}
\end{equation}

Comparing with Eq.~\eqref{eq:mismatch_energy}:
\begin{equation}
    \Kpin = \bsigma \times A_{\text{contact}}
\end{equation}

\subsection{The Geometric Factor $f$}

Not all of the contact surface contributes to mismatch energy. The mismatch
stress penetrates only a thin boundary layer, not the full contact area.

\begin{derivationbox}[title=Effective Area and $f$ Factor]
    Define the effective area fraction $f$:
    \begin{equation}
        A_{\text{eff}} = f \times A_{\text{contact}}
    \end{equation}

    The factor $f$ is the ratio of the stress penetration depth $w_{\text{eff}}$
    to the contact radius $r_{\text{contact}}$:
    \begin{align}
        r_{\text{contact}} &:= \sqrt{\delta \cdot L_0}
            && \text{(from Eq.~\eqref{eq:contact_area})} \\[0.5em]
        w_{\text{eff}} &:= \delta
            && \text{(penetration depth $\sim$ brane thickness)} \\[0.5em]
        f &:= \frac{w_{\text{eff}}}{r_{\text{contact}}}
            = \frac{\delta}{\sqrt{\delta \cdot L_0}}
            = \sqrt{\frac{\delta}{L_0}}
            \approx 0.32
        \label{eq:f_factor}
    \end{align}
    \tagI{} (identified ansatz; derivation remains OPEN)
\end{derivationbox}

\textbf{Physical interpretation:} The factor $f = \sqrt{\delta/L_0}$ is a
penetration-depth ratio. Mismatch stress affects only a layer of thickness
$\delta$ at the contact boundary, while the full contact radius is
$\sqrt{\delta L_0} \approx 0.32\;\text{fm}$. This ansatz is geometrically
motivated but not rigorously derived from the 5D action.

% ============================================================================
\section{The $f$ Factor: Status and Open Questions}
\label{sec:sigmaK:f_open}
% ============================================================================

The factor $f = \sqrt{\delta/L_0} \approx 0.32$ is \textbf{identified} [I], not
\textbf{derived} [Der].

\begin{edcopenbox}
    \textbf{OPEN Problem 4.1} \tagOPEN \\
    Derive the geometric factor $f$ from first principles using:
    \begin{enumerate}
        \item Detailed 5D geometry of junction contact
        \item Calculation of curvature-induced stress
        \item Integration over contact surface with bulk+brane+GHY+Israel action
    \end{enumerate}
    Current status: $f = \sqrt{\delta/L_0}$ is a reasonable ansatz based on
    dimensional analysis, but a rigorous derivation is needed.
\end{edcopenbox}

% ============================================================================
\section{The Pinning Constant: Final Formula}
\label{sec:sigmaK:final}
% ============================================================================

\begin{derivationbox}[title=Pinning Constant Derivation]
    Combining all factors:
    \begin{equation}
        \boxed{\Kpin = f \times \bsigma \times A_{\text{contact}}} \tagDcI
        \label{eq:K_formula}
    \end{equation}
    With:
    \begin{align}
        \bsigma &= 8.82\;\MeV/\text{fm}^2 & \tagDc \\
        A_{\text{contact}} &= \pi\delta L_0 = 0.33\;\text{fm}^2 & \tagDc \\
        f &= \sqrt{\delta/L_0} = 0.32 & \tagI
    \end{align}
\end{derivationbox}

% ============================================================================
\section{Numerical Value}
\label{sec:sigmaK:numerical}
% ============================================================================

\subsection{Model Calculation}

\begin{equation}
    \Kpin = 0.32 \times 8.82 \times 0.33 \approx 0.93\;\MeV \tagDcI
    \label{eq:K_numerical}
\end{equation}

\subsection{Consistency Check (Not a Fit)}

\begin{tcolorbox}[colback=yellow!5,colframe=orange!50!black,title=Anti-Tuning Declaration]
The value $\Kpin \approx 0.93\;\MeV$ computed above is \emph{not} adjusted to
match observation. The parameters $\bsigma$, $\delta$, $L_0$, and $f$ are fixed
by their respective sources (Book~I, Ch.~\ref{ch:L0delta}, and dimensional
ansatz). This subsection compares the model output to an independent benchmark;
\textbf{no parameter is tuned here}.
\end{tcolorbox}

\textbf{Observer-level benchmark:} Chapters~\ref{ch:deuterium}--\ref{ch:helium4}
derive cluster binding energies using $\Kpin$. Working backwards from measured
binding energies (treated as baselines [BL]):
\begin{itemize}
    \item From deuterium binding $B_d = 2.22\;\MeV$: implies $\Kpin \approx 0.74\;\MeV$
    \item From closed-4 energy budget: implies $\Kpin \approx 0.8\;\MeV$
\end{itemize}

\textbf{Comparison:}
\begin{center}
\begin{tabular}{lll}
    \toprule
    Source & $\Kpin$ (MeV) & Status \\
    \midrule
    Model (this chapter) & 0.93 & [Dc/I] \\
    Benchmark (from binding) & 0.74--0.80 & [Cal] (Appendix~Q) \\
    \bottomrule
\end{tabular}
\end{center}

The model overshoots the benchmark by $\sim 15$--$25\%$. This discrepancy is
\emph{not} hidden or tuned away. Possible sources:
\begin{itemize}
    \item The $f$-factor ansatz $f = \sqrt{\delta/L_0}$ may be too large
    \item Higher-order corrections to the curvature energy are neglected
    \item The contact geometry model (saddle surface) is approximate
\end{itemize}

Closing this gap requires a rigorous derivation of $f$ from the 5D action
(OPEN Problem~4.1). Until then, the discrepancy is an honest limitation of
the current model.

\subsection{Sensitivity Analysis}

\begin{center}
\begin{tabular}{ccc}
    \toprule
    $f$ value & $\Kpin$ (MeV) & Status \\
    \midrule
    0.25 & 0.73 & Lower bound \\
    0.32 & 0.93 & Central (model) \\
    0.35 & 1.02 & Upper bound \\
    \bottomrule
\end{tabular}
\end{center}

The range $\Kpin \in [0.7, 1.0]\;\MeV$ is robust against $\pm 20\%$ variation in $f$.

% ============================================================================
\section{Input/Output Summary}
\label{sec:sigmaK:table}
% ============================================================================

\begin{table}[htbp]
\centering
\caption{Derivation Chain: $\bsigma \to \Kpin$}
\label{tab:sigmaK:chain}
\begin{tabular}{llll}
    \toprule
    Quantity & Value & Tag & Source \\
    \midrule
    \multicolumn{4}{l}{\textit{Inputs}} \\
    Brane tension $\bsigma$ & $8.82\;\MeV/\text{fm}^2$ & [Dc] & Ch.~\ref{ch:sigma_to_K} \\
    Junction scale $L_0$ & $1\;\text{fm}$ & [P] & Ch.~\ref{ch:L0delta} \\
    Junction width $\delta$ & $0.105\;\text{fm}$ & [P] & Ch.~\ref{ch:L0delta} \\
    \midrule
    \multicolumn{4}{l}{\textit{Derived}} \\
    Contact area $A_{\text{contact}}$ & $0.33\;\text{fm}^2$ & [Dc] & Eq.~\eqref{eq:contact_area} \\
    Geometric factor $f$ & $0.32$ & [I] & Eq.~\eqref{eq:f_factor} \\
    \midrule
    \multicolumn{4}{l}{\textit{Output}} \\
    Pinning constant $\Kpin$ & $0.93\;\MeV$ & [Dc/I] & Eq.~\eqref{eq:K_numerical} \\
    Phenomenological $\Kpin$ & $0.74\;\MeV$ & [Cal] & App.~\ref{app:quarantine} \\
    \bottomrule
\end{tabular}
\end{table}

% ============================================================================
% Observer-side projection
% ============================================================================

\begin{observerbox}
Observer-side projection note: quoted terms are measurement labels for the
3D/4D projection of the 5D objects defined in this chapter; they do not
imply any conventional mechanism.

\begin{itemize}
    \item \ClusterState{} $\leftrightarrow$ measured bound system (projection label)
    \item Pinning energy $\Kpin$ $\leftrightarrow$ binding energy scale (projection label)
\end{itemize}

What the observer records: the pinning constant $\Kpin$ derived from brane
tension $\bsigma$ determines the binding energy scale measured in cluster
systems at the projection level. The observer labels this as ``binding
energy''; the 5D origin is topological pinning.
\end{observerbox}

% ============================================================================
\section{Summary}
\label{sec:sigmaK:summary}
% ============================================================================

\begin{resultbox}
\textbf{Chapter 4 Result:}
\begin{itemize}
    \item Pinning constant formula: $\Kpin = f \times \bsigma \times A_{\text{contact}}$
    \item Contact area: $A_{\text{contact}} = \pi\delta L_0 \approx 0.33\;\text{fm}^2$ [Dc]
    \item Geometric factor: $f = \sqrt{\delta/L_0} \approx 0.32$ [I]
    \item Model result: $\Kpin \approx 0.93\;\MeV$ [Dc/I]
    \item Phenomenological: $\Kpin \approx 0.74\;\MeV$ [Cal]
    \item Agreement: Within 15--20\%
\end{itemize}
All cluster binding traces to the single fundamental parameter $\bsigma$.
\end{resultbox}

\begin{edcopenbox}
\textbf{OPEN Problem 4.2} \tagOPEN \\
Close the 15--20\% gap between model $\Kpin = 0.93\;\MeV$ and phenomenological
$\Kpin = 0.74\;\MeV$ by:
\begin{enumerate}
    \item Rigorous derivation of $f$ factor
    \item Refinement of contact geometry model
    \item Higher-order corrections to curvature energy
\end{enumerate}
\end{edcopenbox}

% ----------------------------------------------------------------------------
% BRIDGE TO NEXT CHAPTER
% ----------------------------------------------------------------------------
\paragraph{Bridge to Chapter~\ref{ch:M6}.}
The pinning constant $\Kpin$ sets the energy scale for individual bonds.
To compute cluster binding energies, we need to count how many bonds form
in a given configuration---this requires the full $\Msix$ coordination
lattice developed in Chapter~\ref{ch:M6}. That chapter derives the
allowed coordination set $S = \{2^a \times 3^b\}$ and identifies
the forbidden zones where frustration effects dominate.

